%% ============================================================
%% EXEMPLO FICTÍCIO — o nível que a equipe espera
%%
%% Tudo aqui é inventado: a pessoa, o professor, os códigos, os números,
%% os resultados. O exemplo existe para mostrar como fica um relatório
%% completo, sem pendência: cada atividade com a ficha, o desenvolvimento
%% técnico com figuras e equações, os resultados medidos, as evidências.
%% Não copie o texto; copie o padrão.
%%
%% As figuras foram desenhadas em TikZ só para o exemplo caber no
%% repositório; no seu relatório, exporte do KiCad, do CAD, do simulador,
%% e use \includegraphics.
%% ============================================================
\documentclass[final,exemplo]{nro-aacc}
\usepackage{pgfplots}
\pgfplotsset{compat=1.18}
\usepackage{circuitikz}
\usetikzlibrary{arrows.meta}

\estudante{
  nome      = Ana Clara Exemplo,
  matricula = 2021000000,
  curso     = Engenharia Elétrica,
  colegiado = Colegiado do Curso de Graduação em Engenharia Elétrica,
  email     = ana.exemplo@ufmg.br,
  ingresso  = 2021/1,
  periodo   = 10º,
}
\membro{
  registro     = 999,
  situacao     = egresso,
  inicio       = 2023-03,
  fim          = 2025-12,
  cargo        = {Supervisora do subsistema eletrônico do NEBULA},
  departamento = {Pesquisa e Desenvolvimento},
  projetos     = NEBULA,
  supervisor   = {Rafael Fictício, supervisor do projeto NEBULA},
  dedicacao    = 10,
  siex         = não,
  vinculo      = NRO-DIR-004-999,
  verificador  = EXEM-PLOF-ICTI,
}
\documento{
  pn   = 999,
  data = 2026-02-10,
}
%% no exemplo, também a equipe é fictícia
\equipe{
  professor             = {Prof. Dr. João Modelo},
  departamentoprofessor = {Departamento de Engenharia Mecânica},
  dirigente             = {Beatriz Amostra},
  cargodirigente        = {Diretora-geral},
  siextitulo            = {Tecnologias assistivas para a reabilitação e o esporte adaptado (exemplo)},
  siexnumero            = {000000},
  siexcoordenador       = {Prof. Dr. João Modelo},
  siexvigencia          = {2024 a 2026},
}

%% o pedido: os códigos e os dispositivos também são de exemplo
\modalidade{COMP}{codigo = {ELE-AC-04 (exemplo)}, dispositivo = {art. X, inciso Y (exemplo)}, limite = 180}
\modalidade{PES}{codigo = {ELE-AC-01 (exemplo)}, dispositivo = {art. X, inciso Y (exemplo)}}
\modalidade{EXT}{codigo = {ELE-AC-02 (exemplo)}, dispositivo = {art. X, inciso Y (exemplo)}}
\modalidade{ENS}{codigo = {ELE-AC-06 (exemplo)}, dispositivo = {art. X, inciso Y (exemplo)}}
\modalidade{EVT}{codigo = {ELE-AC-03 (exemplo)}, dispositivo = {art. X, inciso Y (exemplo)}}
\modalidade{GES}{codigo = {ELE-AC-07 (exemplo)}, dispositivo = {art. X, inciso Y (exemplo)}}

\addbibresource{exemplo/referencias-exemplo.bib}

\begin{document}

\capa
\requerimento
\quadrospedido
\sumario

\parteinstitucional

%% ============================================================
\parte[Quem é a estudante, como ela chegou à equipe e o que fez nela, em linhas gerais. As atividades,
  uma a uma, estão na Parte~III.]{A estudante na equipe}

\identificacao

\section{Resumo da participação}
Participei da NeuroDynamics por quase três anos, de março de 2023 a dezembro de 2025, sempre no
subsistema eletrônico do NEBULA, o estimulador elétrico funcional que leva uma pessoa com lesão medular
a pedalar. Fui a responsável pela revisão B da placa de potência, do esquemático à placa montada e
testada, e pela simulação que definiu a tensão de alimentação dela; projetei o invólucro isolado;
conduzi os ensaios de bancada e a parte técnica das sessões de validação com o atleta-piloto; operei o
estimulador numa competição internacional; escrevi o treinamento de segurança elétrica que os trainees
fazem hoje; e, em 2025, supervisionei o subsistema até a liberação da revisão C.

É trabalho de Engenharia Elétrica do começo ao fim --- eletrônica analógica e de potência,
compatibilidade eletromagnética, instrumentação, controle ---, com um usuário real, com as normas de
equipamento eletromédico como referência e com a documentação que um produto exige. Peço 440 horas,
e só as que têm evidência.

\section{Trajetória na equipe}
\begin{trajetoria}
  \marco{2023-03}{Aprovação no processo seletivo de 2023/1 (dinâmica em grupo e entrevista) e início do período trainee.}
  \marco{2023-06}{Apresentação do desafio trainee e entrada no NEBULA, no subsistema eletrônico.}
  \marco{2023-12}{Afastamento temporário de dois meses, aprovado pelo Depto. de Pessoal (estágio de férias).}
  \marco{2024-03}{Responsável pela revisão B da placa de potência (Atividades 1 a 3).}
  \marco{2024-10}{Técnica do estimulador na competição internacional (Atividade 6).}
  \marco{2025-01}{Progressão a supervisora do subsistema eletrônico (Atividade 7).}
  \marco{2025-03}{Autora e instrutora do treinamento de segurança elétrica (Atividade 5).}
  \marco{2025-12}{Desligamento a pedido, pela conclusão do curso, com a transição do subsistema ao sucessor.}
\end{trajetoria}

\section{Afastamentos e interrupções}
\begin{afastamentos}
  \afastamento{2023-12}{2024-01}{afastamento temporário aprovado pelo Depto. de Pessoal (\texttt{NRO-PES-002-990}), para um estágio de férias.}
\end{afastamentos}

\section{Formação recebida na equipe}
\begin{treinamentos}
  \treinamento{NRO-TRE-901}{Integração à NeuroDynamics}{3}{2023-04-02}{CERT-EXEM-0001}
  \treinamento{NRO-TRE-902}{Controle de documentos e registros}{2}{2023-05-20}{CERT-EXEM-0002}
  \treinamento{NRO-TRE-903}{Segurança elétrica em bancada (versão de 2022)}{2}{2023-05-10}{CERT-EXEM-0003}
  \treinamento{NRO-TRE-904}{Boas práticas de projeto de placas}{6}{2023-09-15}{CERT-EXEM-0004}
\end{treinamentos}

%% ============================================================
\parteatividades

%% EXEMPLO FICTÍCIO — Atividade 1: eletrônica e placa de circuito impresso
\begin{atividade}{Placa de potência do estimulador do NEBULA, revisão B}{
  tipo         = pcb,
  modalidade   = COMP,
  alternativa  = PES,
  horas        = 150,
  inicio       = 2024-03,
  fim          = 2024-09,
  projeto      = {NEBULA, estimulação elétrica funcional para ciclismo},
  papel        = Responsável,
  supervisor   = {Rafael Fictício, supervisor do projeto NEBULA},
  registros    = {NBL-114; NBL-121; NBL-133; NRO-PRO-014-901; NRO-PRO-003-901},
  competencias = {I, II, III, V},
  disciplinas  = {Eletrônica Analógica; Eletrônica de Potência; Circuitos Elétricos II; Compatibilidade Eletromagnética},
}

\begin{contexto}
O NEBULA é o projeto da equipe que leva uma pessoa com lesão medular a pedalar uma bicicleta
adaptada por meio de estimulação elétrica funcional (FES): pulsos de corrente, aplicados por eletrodos
de superfície, contraem na ordem certa os músculos das pernas. O coração do sistema é o estimulador,
e a placa de potência é a parte dele que gera os pulsos.

A revisão A da placa funcionava, mas os testes de 2023 (\texttt{NRO-PRO-003-900}) registraram quatro
problemas: a amplitude da corrente caía quando a impedância do eletrodo subia; o estágio de saída
chegava a \SI{71}{\celsius} em estimulação contínua; o chaveamento da saída induzia ruído no sinal do
encoder do pedal; e não havia limite de corrente por hardware --- a proteção dependia só do firmware.

Os requisitos vêm da USRS do projeto (\texttt{NRO-PRO-004-901}): pulsos bifásicos com carga
balanceada; amplitude de 0 a \SI{100}{\milli\ampere}, em passos de \SI{1}{\milli\ampere}; largura de
50 a \SI{500}{\micro\second}; frequência de 20 a \SI{50}{\hertz}; quatro canais; e isolação galvânica
entre o usuário e todo ponto ligado à rede elétrica (o carregador e o USB). A revisão B tinha de
resolver os quatro problemas sem sair desses requisitos.
\end{contexto}

\begin{contribuicao}
Fui a responsável pela revisão B, do esquemático à placa montada e testada:
\begin{itemize}
  \item reprojetei o estágio de saída (a fonte de corrente e a ponte H) e criei o limite de corrente
    por hardware, que não existia;
  \item escolhi os componentes críticos --- o amplificador operacional, os MOSFETs de \SI{250}{\volt}, os
    isoladores --- e registrei as escolhas e as alternativas no design record (\texttt{NRO-PRO-014-901});
  \item fiz o layout das quatro camadas, montei as duas primeiras placas no LABBIO e conduzi a primeira
    energização.
\end{itemize}
Não foram meus: o firmware que gera os pulsos (Lucas Fictício); o conversor boost, que veio da
revisão A e em que só ajustei a tensão de saída; e a revisão do layout antes da fabricação, feita pelo
supervisor do projeto.
\end{contribuicao}

\begin{desenvolvimento}
\fichatecnica{
  ferramenta  = {KiCad 8.0},
  revisao     = {B (a revisão C incorpora a errata desta)},
  camadas     = {4 camadas, FR-4 de 1,6\,mm},
  dimensoes   = {100 × 80\,mm},
  componentes = {146},
  alimentacao = {Bateria 3S (11,1\,V); DC-DC isolado de 3\,kV para o lado do usuário; boost de 120\,V},
  fabricacao  = {Fabricação terceirizada; montagem manual no LABBIO (duas placas)},
  normas      = {IEC 60601-1 e IEC 60601-2-10 como referência de projeto, sem certificação},
  arquivos    = {Repositório \texttt{nebula-hw}, tag \texttt{rev-b}},
}

\textbf{A arquitetura.} A placa tem dois lados, separados por uma barreira de isolação
(Figura~\ref{fig:ex-blocos}). Do lado ligado ao USB e ao carregador ficam o microcontrolador e a
bateria; do lado do usuário, o conversor boost que gera os \SI{120}{\volt}, os conversores
digital-analógicos que definem a corrente de cada canal e os estágios de saída. Só atravessam a barreira
a energia, por um conversor DC-DC isolado, e os sinais digitais, por isoladores. A revisão A não tinha
essa separação: o USB e o eletrodo dividiam o mesmo terra.

\begin{figure}[htbp]
  \centering
  \resizebox{\linewidth}{!}{% Diagrama de blocos da placa de potência, revisão B (exemplo fictício)
\begin{tikzpicture}[
    bloco/.style={draw=nro-tinta, line width=0.5pt, fill=white, rounded corners=1.2pt,
      minimum height=11mm, text width=29mm, align=center, font=\footnotesize, inner sep=1.5pt},
    critico/.style={bloco, fill=nro-synapse!18},
    seta/.style={-{Stealth[length=1.8mm]}, line width=0.55pt, draw=nro-tinta},
    alta/.style={seta, draw=nro-pulso-texto, line width=0.8pt},
    x=1mm, y=1mm]
  % colunas: controle | barreira | lado do usuário
  \node[bloco]   (mcu)   at (0,32)  {Microcontrolador\\{\scriptsize firmware, USB}};
  \node[bloco]   (bat)   at (0,0)   {Bateria 3S\\{\scriptsize 11,1\,V, 2,6\,Ah}};
  \node[bloco]   (iso)   at (44,32) {Isoladores digitais\\{\scriptsize SPI e habilitações}};
  \node[bloco]   (dcdc)  at (44,0)  {DC-DC isolado\\{\scriptsize 3\,kV, 12\,V}};
  \node[bloco]   (dac)   at (88,32) {DAC de 12 bits\\{\scriptsize a referência de corrente}};
  \node[bloco]   (boost) at (88,0)  {Conversor boost\\{\scriptsize 12\,V para 120\,V}};
  \node[critico] (fonte) at (132,32) {Fonte de corrente\\{\scriptsize amp. op. e MOSFET}};
  \node[critico] (ponte) at (132,0)  {Ponte H\\{\scriptsize pulso bifásico}};
  \node[critico] (lim)   at (88,16) {\scriptsize Limite de corrente: 110\,mA, trava e parada de emergência};
  \node[bloco, text width=22mm] (ele) at (132,-17) {Eletrodos\\{\scriptsize 4 canais}};
  \draw[seta] (mcu) -- (iso);
  \draw[seta] (iso) -- (dac);
  \draw[seta] (dac) -- (fonte);
  \draw[seta] (bat) -- (dcdc);
  \draw[seta] (dcdc) -- (boost);
  \draw[alta] (boost) -- node[above, font=\scriptsize, text=nro-pulso-texto] {120\,V} (ponte);
  \draw[seta] (fonte) -- (ponte);
  \draw[seta] (ponte) -- (ele);
  \draw[seta] ([xshift=-10mm]fonte.south) |- ([yshift=2mm]lim.east);
  \draw[seta] ([yshift=-2.5mm]lim.east) -| ([xshift=-7mm]ponte.north);
  \draw[seta] (bat.north) -- (mcu.south);
  % a barreira de isolação
  \draw[dashed, line width=0.7pt, nro-pulso-texto] (44,-9) -- (44,43)
    node[above, font=\scriptsize, text=nro-pulso-texto] {barreira de isolação (3\,kV)};
  \node[font=\scriptsize, text=nro-cinza, anchor=west] at (-14,41) {lado ligado ao USB e ao carregador};
  \node[font=\scriptsize, text=nro-cinza, anchor=east] at (146,41) {lado do usuário};
\end{tikzpicture}
}
  \caption{Diagrama de blocos da placa de potência, revisão B. Em verde-claro, os blocos refeitos nesta
    revisão.}
  \fonte{Elaborado pela autora.}
  \label{fig:ex-blocos}
\end{figure}

\textbf{A fonte de corrente.} Cada canal tem uma fonte de corrente controlada por tensão
(Figura~\ref{fig:ex-fonte}): o amplificador operacional força sobre o resistor de medição $R_S$ a
tensão que o DAC define, e a mesma corrente atravessa a carga. Com ganho de malha alto,
\begin{equation}
  I_\text{saída} = \frac{V_\mathrm{DAC}}{R_S}.
  \label{eq:ex-fonte}
\end{equation}
Com $R_S = \SI{10}{\ohm}$, a faixa de 0 a \SI{100}{\milli\ampere} pede $V_\mathrm{DAC}$ de 0 a
\SI{1}{\volt}; o DAC de 12 bits dá passos de \SI{24}{\micro\ampere}, bem abaixo do passo de
\SI{1}{\milli\ampere} do requisito. Escolhi um resistor de 0,1\,\% e 25\,ppm/°C: a exatidão da
corrente é a dele.

A queda de amplitude da revisão A tinha duas causas, que a simulação da Atividade~2 separou: a
tensão de \SI{90}{\volt} não bastava para manter \SI{100}{\milli\ampere} em eletrodos de impedância
alta, e o laço oscilava na borda do pulso. A revisão B sobe a tensão para \SI{120}{\volt} e tem o
capacitor de compensação $C_c = \SI{22}{\pico\farad}$, que a simulação dimensionou.

\begin{figure}[htbp]
  \centering
  % A fonte de corrente de um canal (exemplo fictício): o amplificador
% operacional força sobre R_S a tensão do DAC; a corrente na carga é V_DAC/R_S.
\begin{circuitikz}[font=\footnotesize, line width=0.5pt]
  \draw (0,0) node[op amp, yscale=-1] (oa) {};
  \node at ([xshift=-1.5mm]oa.center) {\scriptsize U1};
  % a referência na entrada não inversora (em cima)
  \draw (oa.+) to[short, -o] ++(-1.4,0) node[left] {$V_\mathrm{DAC}$ (0 a 1\,V)};
  % a saída, o resistor de gate e o MOSFET
  \draw (oa.out) to[R, l=$R_G$, -*] ++(2.4,0) coordinate (g);
  \draw (g) node[nigfete, anchor=G] (q) {};
  \node[right, font=\scriptsize] at (q.east) {Q1 (250\,V)};
  \draw (q.D) -- ++(0,0.5) to[R, l_=$Z_\mathrm{carga}$] ++(0,2.1) node[above] {$V_\mathrm{HV} = 120$\,V};
  % o resistor de medição e a realimentação, por baixo
  \draw (q.S) -- ++(0,-1.9) coordinate (s) to[R, l_={$R_S = 10\,\Omega$}, *-] ++(0,-2) node[ground] {};
  \coordinate (r0) at ([xshift=-9mm]oa.-);
  \draw (s) -- (s -| r0) |- (oa.-);
  \draw (s) -- ++(1.4,0) node[right, align=left, font=\scriptsize] {ao comparador\\do limite (110\,mA)};
  % a compensação, entre a saída e a entrada inversora
  \draw ([xshift=-4mm]oa.-) node[circ] {} -- ++(0,-0.9) coordinate (c1);
  \draw (c1) to[C, l_=$C_c$] (c1 -| oa.out) -- (oa.out);
\end{circuitikz}

  \caption{A fonte de corrente de um canal. A tensão em $R_S$ também alimenta o comparador do limite de
    corrente.}
  \fonte{Elaborado pela autora.}
  \label{fig:ex-fonte}
\end{figure}

\textbf{A dissipação.} No pior caso --- carga em curto ---, o MOSFET dissipa toda a tensão:
\begin{equation}
  P_\mathrm{pico} = V_\mathrm{HV}\, I = \SI{120}{\volt} \times \SI{0,1}{\ampere} = \SI{12}{\watt},
  \qquad
  P_\text{média} = P_\mathrm{pico}\, (2\, t_p\, f) = \SI{12}{\watt} \times 0{,}03 = \SI{0,36}{\watt},
  \label{eq:ex-potencia}
\end{equation}
com $t_p = \SI{300}{\micro\second}$ e $f = \SI{50}{\hertz}$. A câmera térmica mostrou que o
aquecimento da revisão A não vinha dos MOSFETs, mas do regulador linear que alimentava os drivers da
ponte, dissipando \SI{1,4}{\watt}. Troquei-o por um regulador chaveado e pus vias térmicas sob cada
MOSFET.

\textbf{O limite de corrente.} Um comparador lê a tensão em $R_S$; acima de \SI{1,1}{\volt}
(\SI{110}{\milli\ampere}), ele aciona uma trava que desliga os drivers da ponte H, sem passar pelo
firmware. A trava só se desfaz por comando do firmware depois de a pessoa que opera apertar o reset,
e o botão de parada de emergência corta a habilitação do boost diretamente.

\textbf{O layout.} Quatro camadas: sinais e componentes em cima, terra contínuo na segunda, as
alimentações de \SI{12}{\volt} e de \SI{120}{\volt} em ilhas separadas na terceira, sinais embaixo
(Figura~\ref{fig:ex-layout}). Entre os dois lados, um rasgo de \SI{2,5}{\milli\meter} e
\SI{8}{\milli\meter} de distância de escoamento, com a tabela da IEC 60601-1 como referência. A malha de
corrente de cada canal ficou tão pequena quanto possível, com ligação Kelvin em $R_S$; o nó chaveado
da ponte ficou longe da entrada do encoder, que ganhou um filtro RC e um Schmitt trigger. A verificação
de regras (DRC) usou \SI{0,3}{\milli\meter} de isolação geral e \SI{1,5}{\milli\meter} na área de alta
tensão, e a verificação elétrica (ERC) terminou sem erro.

\begin{figure}[htbp]
  \centering
  % Planta simplificada do layout da revisão B (exemplo fictício): as áreas,
% a barreira de isolação e os estágios de saída. Não é o layout completo.
\begin{tikzpicture}[x=0.95mm, y=0.95mm, font=\scriptsize]
  % a placa, 100 x 80 mm
  \draw[line width=0.8pt, fill=nro-axon!8, rounded corners=1.5] (0,0) rectangle (100,80);
  \foreach \x/\y in {4/4,96/4,4/76,96/76} \draw[fill=white] (\x,\y) circle (1.6);
  % lado ligado ao USB
  \draw[densely dashed, nro-cinza, fill=nro-fundoclaro] (7,8) rectangle (33,72);
  \node[align=center, text=nro-cinza] at (20,67) {controle};
  \draw[fill=nro-cinza!45] (12,42) rectangle (26,56) node[midway, text=white] {MCU};
  \draw[fill=nro-cinza!25] (8,18) rectangle (16,24) node[midway, font=\tiny] {USB};
  \draw[fill=nro-cinza!25] (18,12) rectangle (30,22) node[midway, font=\tiny, align=center] {carga da\\bateria};
  % a barreira: o rasgo e os componentes que a cruzam
  \fill[white] (36,6) rectangle (38.5,74);
  \draw[nro-pulso-texto, line width=0.4pt] (36,6) rectangle (38.5,74);
  \draw[fill=nro-synapse!45] (32,28) rectangle (43,38) node[midway, font=\tiny, align=center] {DC-DC\\isolado};
  \draw[fill=nro-synapse!45] (32,48) rectangle (43,56) node[midway, font=\tiny, align=center] {isol.\\digitais};
  % lado do usuário (alta tensão)
  \draw[densely dashed, nro-pulso-texto, fill=nro-pulso!7] (46,8) rectangle (93,72);
  \node[text=nro-pulso-texto] at (69.5,67) {lado do usuário (120\,V)};
  \draw[fill=nro-tinta!55] (49,48) rectangle (63,60) node[midway, text=white, align=center, font=\tiny] {boost\\120\,V};
  \draw[fill=nro-tinta!30] (49,30) rectangle (61,40) node[midway, font=\tiny, align=center] {DAC e\\fontes};
  % os quatro estágios de saída, com as vias térmicas
  \foreach \i in {0,...,3} {
    \draw[fill=nro-tinta!75] (66+\i*6.5,14) rectangle (71+\i*6.5,22);
    \foreach \a in {0,1,2} \fill[nro-synapse!80!black] (66.8+\i*6.5+\a*1.6,12.2) circle (0.4);
  }
  \node[align=center, font=\tiny] at (79,27) {estágios de saída (Q1--Q4)\\e vias térmicas};
  % as saídas
  \foreach \i in {0,...,3} \draw[fill=nro-synapse!70] (95,14+\i*14) rectangle (100,22+\i*14);
  \node[anchor=west, align=left] at (101.5,40) {4 saídas\\(eletrodos)};
  % cotas
  \draw[{Stealth[length=1.4mm]}-{Stealth[length=1.4mm]}] (0,-4) -- (100,-4) node[midway, below] {100\,mm};
  \draw[{Stealth[length=1.4mm]}-{Stealth[length=1.4mm]}] (-4,0) -- (-4,80) node[midway, left] {80\,mm};
  \draw[nro-pulso-texto] (37.25,74) -- (37.25,86) node[above, align=center] {rasgo de 2,5\,mm: 8\,mm de\\distância de escoamento};
\end{tikzpicture}

  \caption{Planta simplificada do layout da revisão B: as áreas, a barreira de isolação e os estágios
    de saída. O layout completo está no repositório.}
  \fonte{Elaborado pela autora.}
  \label{fig:ex-layout}
\end{figure}
\end{desenvolvimento}

\begin{resultados}
As duas placas foram ensaiadas na bancada com o procedimento do plano de validação; a Atividade~4
descreve a montagem. A Tabela~\ref{tab:ex-revb} compara as revisões com os requisitos.

\begin{table}[htbp]
  \centering\small
  \caption{Revisão A e revisão B, contra os requisitos (média de 5 medições, placa 1).}
  \label{tab:ex-revb}
  \begin{tabular}{|l|c|c|c|}\hline
    \rowcolor{nro-fundo}\rotulo{Requisito} & \rotulo{Esperado} & \rotulo{Rev. A} & \rotulo{Rev. B} \\ \hline
    Amplitude em \SI{1}{\kilo\ohm}, \SI{100}{\milli\ampere} programados & $100 \pm 2$\,mA & 84\,mA & 100,9\,mA \\ \hline
    Amplitude em \SI{500}{\ohm}, \SI{100}{\milli\ampere} programados & $100 \pm 2$\,mA & 99\,mA & 100,4\,mA \\ \hline
    Temperatura do estágio de saída (60\,min, 50\,Hz) & $\le \SI{60}{\celsius}$ & \SI{71}{\celsius} & \SI{48}{\celsius} \\ \hline
    Ruído no sinal do encoder (pico a pico) & $\le \SI{50}{\milli\volt}$ & \SI{420}{\milli\volt} & \SI{35}{\milli\volt} \\ \hline
    Disparo do limite de corrente & 105 a 115\,mA & não havia & 108\,mA, em \SI{1,6}{\micro\second} \\ \hline
  \end{tabular}
  \fonte{Relatório de execução \texttt{NRO-PRO-003-901}.}
\end{table}

A primeira energização revelou dois defeitos, ambos meus: faltava o resistor de pull-down na
habilitação da trava, e as saídas podiam pulsar na partida; e o footprint do DC-DC isolado estava
espelhado. Os dois foram corrigidos à mão nas duas placas e estão entre os sete itens da errata que
virou a revisão C (\texttt{NBL-140}).
\end{resultados}

\begin{evidencias}
\evidencia{Repositório}{nebula-hw, tag rev-b}{O esquemático, o layout e a lista de materiais da revisão B, com os meus commits}{GitHub da equipe (acesso a pedido)}
\evidencia{Arquivo controlado}{NRO-PRO-014-901}{O design record da revisão B: requisitos, escolhas e alternativas}{SOMA › Arquivos}
\evidencia{Arquivo controlado}{NRO-PRO-003-901}{O relatório de execução dos ensaios de bancada da revisão B}{SOMA › Arquivos}
\evidencia{Atividade}{NBL-114; NBL-121; NBL-133}{Os cartões do esquemático, do layout e da montagem, atribuídos a mim e concluídos}{SOMA › Atividades › NEBULA}
\evidencia{Arquivo controlado}{NRO-PRO-004-901}{A USRS, de onde vêm os requisitos elétricos citados}{SOMA › Arquivos}
\end{evidencias}

\begin{aprendizados}
O curso me deu a fonte de corrente com amplificador operacional (Eletrônica Analógica), o
dimensionamento da dissipação (Eletrônica de Potência) e o transitório de primeira ordem (Circuitos
Elétricos II). O que ele não deu, e a placa deu: projetar para a segurança de uma pessoa ligada ao
circuito, ler uma norma de equipamento eletromédico, fazer layout pensando em ruído e em isolação, e
descobrir que a causa de um problema raramente é a primeira suspeita --- o calor não vinha de onde eu
esperava. Na próxima revisão, eu faria a revisão do footprint dos componentes novos contra o datasheet
antes de mandar fabricar, e não depois.
\end{aprendizados}
\end{atividade}

%% EXEMPLO FICTÍCIO — Atividade 2: simulação computacional
\begin{atividade}{Simulação do estágio de saída e do limite de tensão de conformidade}{
  tipo         = simulacao,
  modalidade   = PES,
  alternativa  = VPC,
  horas        = 45,
  inicio       = 2024-02,
  fim          = 2024-04,
  projeto      = {NEBULA, estimulação elétrica funcional para ciclismo},
  papel        = Responsável,
  supervisor   = {Rafael Fictício, supervisor do projeto NEBULA},
  registros    = {NBL-109; NRO-PRO-013-901},
  competencias = {II, III, VIII},
  disciplinas  = {Circuitos Elétricos II; Eletrônica Analógica; Sistemas de Controle; Cálculo Numérico},
}

\begin{contexto}
Antes de refazer o estágio de saída (Atividade~1), a equipe precisava responder a duas perguntas.
Primeira: com a impedância real dos eletrodos, que tensão de alimentação garante
\SI{100}{\milli\ampere} durante um pulso inteiro de \SI{300}{\micro\second}? Segunda: o laço do
amplificador operacional com o MOSFET, que oscilava na revisão A, fica estável com os componentes
novos? Da primeira dependia a escolha entre manter \SI{90}{\volt} e subir para \SI{120}{\volt} ---
com MOSFETs, capacitores e isolação mais caros; da segunda, o capacitor de compensação.
\end{contexto}

\begin{contribuicao}
Montei o modelo da carga a partir da literatura, levantei os parâmetros dos componentes, rodei as
simulações e as varreduras, validei o resultado contra a medição na placa e escrevi o registro de
decisão de arquitetura que fixou a tensão em \SI{120}{\volt} (\texttt{NRO-PRO-013-901}). O modelo do
MOSFET é o do fabricante; o do amplificador operacional, o macromodelo do fabricante, sem alteração.
\end{contribuicao}

\begin{desenvolvimento}
\fichatecnica{
  ferramenta    = {LTspice 24},
  fenomeno      = {A tensão no eletrodo durante o pulso e a resposta da fonte de corrente à borda do pulso},
  modelo        = {Circuito elétrico: fonte de corrente da placa e carga no modelo de Randles; análise transitória e de pequenos sinais},
  parametros    = {$R_s = 200\,\Omega$; $R_p$ de 500\,$\Omega$ a 2\,k$\Omega$; $C_p = 100$\,nF (faixas de Merrill, Bikson e Jefferys, 2005); modelos dos fabricantes para U1 e Q1},
  discretizacao = {Passo máximo de 10\,ns; com 5\,ns, o sobressinal mudou menos de 0,2 ponto percentual},
  validacao     = {Comparação com a placa montada, em carga resistiva de 1\,k$\Omega$ (Atividade~4)},
  arquivos      = {Repositório \texttt{nebula-sim}, pasta \texttt{saida/}},
}

\textbf{O modelo da carga.} O conjunto eletrodo e pele se comporta, em primeira aproximação, como
um resistor em série com um paralelo RC --- o modelo de Randles (Figura~\ref{fig:ex-randles}) ---, com
valores na faixa relatada por \textcite{merrill2005}. Sob um degrau de corrente $I$, a tensão no
eletrodo é
\begin{equation}
  v_e(t) = I \left[ R_s + R_p \left( 1 - e^{-t/(R_p C_p)} \right) \right].
  \label{eq:ex-randles}
\end{equation}
A fonte mantém a corrente enquanto $v_e$ somada às quedas em Q1 e em $R_S$ (cerca de
\SI{2}{\volt}) não passa da tensão de alimentação: com \SI{120}{\volt}, a conformidade é de
$V_c \approx \SI{118}{\volt}$. Resolvendo \eqref{eq:ex-randles} para $v_e = V_c$, o instante em que a
corrente começa a cair é
\begin{equation}
  t_\mathrm{sat} = -R_p C_p \ln\!\left( 1 - \frac{V_c/I - R_s}{R_p} \right),
  \label{eq:ex-tsat}
\end{equation}
que só existe quando $I (R_s + R_p) > V_c$.

\begin{figure}[htbp]
  \centering
  % O modelo da carga (exemplo fictício): a fonte de corrente sobre o
% conjunto eletrodo e pele, no modelo de Randles.
\begin{circuitikz}[font=\footnotesize, line width=0.5pt, scale=0.85, transform shape]
  \draw (0,0) to[I, l=$i(t)$] (0,3) to[R, l=$R_s$, -*] (3,3) coordinate (a);
  \draw (a) -- (3,3.9) to[R, l=$R_p$] (6.2,3.9) -- (6.2,3) coordinate (b);
  \draw (a) -- (3,2.1) to[C, l_=$C_p$] (6.2,2.1) -- (b);
  \draw (b) to[short, *-] (7.6,3) -- (7.6,0) -- (0,0);
  \draw (7.6,0) node[ground] {};
  \draw[nro-cinza, densely dashed, rounded corners=2pt] (1.3,1.5) rectangle (6.9,4.8);
  \node[nro-cinza, font=\scriptsize, anchor=south] at (4.1,4.8) {eletrodo e pele (modelo de Randles)};
  \draw[-{Stealth[length=1.6mm]}, nro-axon] (8.4,0.3) -- node[right, text=nro-axon] {$v_e(t)$} (8.4,2.9);
\end{circuitikz}

  \caption{O modelo da carga: a fonte de corrente sobre o eletrodo e a pele (modelo de Randles).}
  \fonte{Elaborado pela autora, a partir de \textcite{merrill2005}.}
  \label{fig:ex-randles}
\end{figure}

\textbf{A estabilidade.} Na análise de pequenos sinais, o laço de U1 com Q1 tinha
\SI{38}{\degree} de margem de fase: a capacitância de entrada do MOSFET novo, maior que a do antigo,
somada à resistência de gate, punha um polo perto do cruzamento. Com $C_c = \SI{22}{\pico\farad}$
entre a saída e a entrada inversora de U1, a margem subiu para \SI{62}{\degree}. Testei 10, 22 e
\SI{47}{\pico\farad}: \SI{47}{\pico\farad} amortecia mais, mas alongava a subida além do que o pulso
de \SI{50}{\micro\second} tolera.
\end{desenvolvimento}

\begin{resultados}
A Figura~\ref{fig:ex-tensao} responde à primeira pergunta. Com \SI{120}{\volt}, o pulso de
\SI{300}{\micro\second} a \SI{100}{\milli\ampere} se sustenta até $R_p = \SI{1}{\kilo\ohm}$
($v_e = \SI{115}{\volt}$ no fim do pulso); com $R_p = \SI{2}{\kilo\ohm}$, a corrente cai a partir de
\SI{135}{\micro\second}, como prevê \eqref{eq:ex-tsat}. Com \SI{90}{\volt}, o limite seria
$R_p \approx \SI{690}{\ohm}$ --- abaixo dos eletrodos já gastos que a equipe usa. A decisão foi subir
para \SI{120}{\volt} e, para impedâncias maiores, fazer o firmware medir a impedância no começo da
sessão e limitar a amplitude: isso virou requisito na revisão C da USRS.

\begin{figure}[htbp]
  \centering
  % Tensão no eletrodo durante um pulso de 100 mA, pela eq. (2), limitada
% pela tensão de conformidade (exemplo fictício).
\begin{tikzpicture}
  \begin{axis}[
      width=0.86\linewidth, height=60mm,
      xlabel={Tempo desde o início do pulso (\si{\micro\second})},
      ylabel={Tensão no eletrodo (\si{\volt})},
      xmin=0, xmax=300, ymin=0, ymax=135,
      grid=major, grid style={nro-fundo},
      tick label style={font=\scriptsize}, label style={font=\footnotesize},
      legend pos=south east, legend cell align=left,
      legend style={font=\scriptsize, draw=nro-linha, fill=white, fill opacity=0.9, text opacity=1}]
    \addplot[domain=0:300, samples=121, very thick, nro-axon] {min(118, 0.1*(200+500*(1-exp(-x/50))))};
    \addlegendentry{$R_p = 500\,\Omega$}
    \addplot[domain=0:300, samples=121, very thick, nro-tinta, dashed] {min(118, 0.1*(200+1000*(1-exp(-x/100))))};
    \addlegendentry{$R_p = 1\,\mathrm{k}\Omega$}
    \addplot[domain=0:300, samples=121, very thick, nro-pulso-texto] {min(118, 0.1*(200+2000*(1-exp(-x/200))))};
    \addlegendentry{$R_p = 2\,\mathrm{k}\Omega$}
    \addplot[domain=0:300, thin, nro-cinza, densely dotted] {118};
    \addlegendentry{conformidade (118\,V)}
    \draw[nro-pulso-texto] (axis cs:134.7,118) circle (1.6pt);
    \node[font=\scriptsize, text=nro-pulso-texto, anchor=south east] at (axis cs:132,119) {a corrente cai a partir de \SI{135}{\micro\second}};
  \end{axis}
\end{tikzpicture}

  \caption{Tensão no eletrodo durante um pulso de \SI{100}{\milli\ampere}, para três valores de $R_p$,
    com a alimentação de \SI{120}{\volt}.}
  \fonte{Elaborado pela autora (simulação, eq.~\ref{eq:ex-randles}).}
  \label{fig:ex-tensao}
\end{figure}

A Figura~\ref{fig:ex-borda} responde à segunda. Sem $C_c$, a simulação mostra 18\,\% de sobressinal
--- a oscilação que a revisão A tinha; com \SI{22}{\pico\farad}, 3\,\%. Na placa montada, o sobressinal
medido foi de 1,9\,\% e o tempo de subida (10 a 90\,\%), de \SI{1,8}{\micro\second}, contra
\SI{1,7}{\micro\second} na simulação: 6\,\% de diferença. A amplitude simulada em \SI{1}{\kilo\ohm}
(\SI{100,0}{\milli\ampere}) ficou a 0,9\,\% da medida.

\begin{figure}[htbp]
  \centering
  % A borda de subida da corrente: simulação sem e com a compensação, e a
% medição na placa montada (exemplo fictício).
\begin{tikzpicture}
  \begin{axis}[
      width=0.86\linewidth, height=58mm,
      xlabel={Tempo (\si{\micro\second})},
      ylabel={Corrente (\% da programada)},
      xmin=0, xmax=8, ymin=0, ymax=125,
      grid=major, grid style={nro-fundo},
      tick label style={font=\scriptsize}, label style={font=\footnotesize},
      legend pos=south east, legend cell align=left,
      legend style={font=\scriptsize, draw=nro-linha, fill=white, fill opacity=0.9, text opacity=1}]
    \addplot[domain=0:8, samples=161, thick, nro-pulso-texto]
      {100*(1-exp(-0.48*1.35*x)*(cos(deg(1.35*sqrt(1-0.48^2)*x)) + 0.48/sqrt(1-0.48^2)*sin(deg(1.35*sqrt(1-0.48^2)*x))))};
    \addlegendentry{simulada, sem $C_c$ (sobressinal de 18\,\%)}
    \addplot[domain=0:8, samples=161, very thick, nro-axon]
      {100*(1-exp(-0.74*1.35*x)*(cos(deg(1.35*sqrt(1-0.74^2)*x)) + 0.74/sqrt(1-0.74^2)*sin(deg(1.35*sqrt(1-0.74^2)*x))))};
    \addlegendentry{simulada, $C_c = 22$\,pF (3\,\%)}
    \addplot[only marks, mark=*, mark size=1.4pt, nro-tinta]
      coordinates {(0.5,13.8) (1,40.2) (1.5,66.6) (2,83.5) (2.5,95.7) (3,100.8) (4,101.9) (5,101.4) (6,99.2) (7,99.8) (8,98.9)};
    \addlegendentry{medida na placa, $C_c = 22$\,pF}
  \end{axis}
\end{tikzpicture}

  \caption{A borda de subida da corrente: simulada sem e com a compensação, e medida na placa.}
  \fonte{Elaborado pela autora.}
  \label{fig:ex-borda}
\end{figure}

\textbf{As limitações.} O modelo de Randles é linear, e a impedância real do eletrodo muda com a
densidade de corrente, com o suor e com o desgaste do gel; os parâmetros vieram da literatura e de
medições em placa de teste, não no usuário. A simulação serve para dimensionar a alimentação e a
compensação; não permite concluir nada sobre o conforto nem sobre a segurança da pele, que ficam com o
protocolo clínico.
\end{resultados}

\begin{evidencias}
\evidencia{Repositório}{nebula-sim, pasta saida/}{Os arquivos da simulação (.asc), os parâmetros e as varreduras}{GitHub da equipe (acesso a pedido)}
\evidencia{Arquivo controlado}{NRO-PRO-013-901}{O registro da decisão de subir a alimentação para 120\,V, com as alternativas}{SOMA › Arquivos}
\evidencia{Atividade}{NBL-109}{O cartão da simulação, atribuído a mim e concluído}{SOMA › Atividades › NEBULA}
\end{evidencias}

\begin{aprendizados}
A simulação juntou três disciplinas numa pergunta só: o transitório de primeira ordem de Circuitos
Elétricos II está em \eqref{eq:ex-randles}, a margem de fase de Sistemas de Controle decidiu o
capacitor, e o cuidado com o passo de tempo veio de Cálculo Numérico. O que aprendi fora da sala: um
modelo vale o que valem os seus parâmetros, e a simulação só ganhou crédito na equipe quando bateu com
a medição. Hoje eu faria a validação antes de apresentar o resultado, e não depois.
\end{aprendizados}
\end{atividade}

%% EXEMPLO FICTÍCIO — Atividade 3: projeto mecânico
\begin{atividade}{Invólucro isolado do estimulador, em impressão 3D}{
  tipo         = mecanico,
  modalidade   = PES,
  alternativa  = VPC,
  horas        = 35,
  inicio       = 2024-06,
  fim          = 2024-08,
  projeto      = {NEBULA, estimulação elétrica funcional para ciclismo},
  papel        = Responsável,
  supervisor   = {Rafael Fictício, supervisor do projeto NEBULA},
  registros    = {NBL-127; NRO-PRO-014-902; NRO-PRO-003-902},
  competencias = {I, III, IV},
  disciplinas  = {Desenho Técnico; Fenômenos de Transporte; Materiais Elétricos},
}

\begin{contexto}
A placa da revisão A viajava numa caixa plástica improvisada: abrir a tampa expunha a alta tensão,
não havia ventilação, e a caixa, presa ao quadro da bicicleta com abraçadeiras de nylon, soltou-se duas
vezes com a vibração da pedalada. Com a placa nova (Atividade~1), o invólucro passou a ser parte do
projeto, com requisitos próprios, registrados no design record: nenhuma parte de alta tensão acessível
sem ferramenta; temperatura interna de no máximo \SI{60}{\celsius} com \SI{35}{\celsius} de ambiente;
massa de até \SI{400}{\gram} com a placa e a bateria; fixação num tubo de \SI{35}{\milli\meter},
removível em menos de um minuto; acesso aos conectores e ao botão de parada; e fabricação na impressora
do LABBIO.
\end{contexto}

\begin{contribuicao}
Concebi o invólucro, comparei as alternativas, fiz o modelo e os desenhos, calculei a temperatura
interna e a abraçadeira, imprimi as peças, montei o conjunto e fiz os ensaios térmico e de vibração. A
fiação do botão de parada foi feita por Lucas Fictício.
\end{contribuicao}

\begin{desenvolvimento}
\fichatecnica{
  ferramenta    = {FreeCAD 0.21 (modelo e desenhos) e PrusaSlicer 2.7 (fatiamento)},
  materiais     = {PETG (base, tampa e abraçadeira); TPU 95A (calço); insertos de latão M3},
  fabricacao    = {Impressão 3D por extrusão (FDM) no LABBIO: camada de 0,2\,mm, 4 perímetros, 30\,\% de preenchimento},
  pecas         = {Base, tampa, abraçadeira e calço; 4 parafusos M3},
  carregamentos = {Massa de até 0,38\,kg sob vibração de até 5\,g; 2,1\,W dissipados dentro},
  desenhos      = {Anexos do design record \texttt{NRO-PRO-014-902}},
  arquivos      = {Drive da equipe, pasta NEBULA/mecânica/involucro-rev1},
}

\textbf{As alternativas.} Comparei três soluções numa matriz de decisão com pesos definidos com o
supervisor --- isolação (30\,\%), desempenho térmico (20\,\%), custo (20\,\%), encaixe na placa
(15\,\%) e prazo (15\,\%): uma caixa comercial de ABS usinada (nota 3,4 de 5), uma caixa de alumínio
usinada (3,1: boa para o calor, mas condutora e cara) e uma peça impressa em PETG (4,3). O PETG
ganhou pela isolação, pelo encaixe sob medida e pela temperatura de transição vítrea, perto de
\SI{80}{\celsius}, bem acima da do PLA.

\textbf{O modelo.} A base recebe a placa sobre quatro colunas com insertos de latão; a tampa tem
ranhuras em labirinto, que deixam o ar passar sem deixar um caminho reto até a alta tensão, e o furo
do botão de parada; a abraçadeira, com calço de TPU, prende o conjunto ao tubo
(Figura~\ref{fig:ex-involucro}). As paredes têm \SI{2,4}{\milli\meter}, múltiplas da largura de
extrusão, e as folgas de encaixe, \SI{0,3}{\milli\meter}.

\begin{figure}[htbp]
  \centering
  % Vista explodida do invólucro (exemplo fictício), em projeção oblíqua:
% a base com a abraçadeira, a placa e a tampa com as ranhuras.
\begin{tikzpicture}[x={(7.2mm,0mm)}, y={(3.1mm,2.2mm)}, z={(0mm,7.2mm)},
    line join=round, font=\scriptsize,
    aresta/.style={draw=nro-tinta, line width=0.45pt}]
  % uma caixa: canto (x,y,z), largura (x), profundidade (y), altura (z), cores
  \def\caixa#1#2#3#4#5#6#7{%
    \filldraw[aresta, fill=#7!70] (#1+#4,#2,#3) -- (#1+#4,#2+#5,#3) -- (#1+#4,#2+#5,#3+#6) -- (#1+#4,#2,#3+#6) -- cycle;
    \filldraw[aresta, fill=#7!45] (#1,#2,#3) -- (#1+#4,#2,#3) -- (#1+#4,#2,#3+#6) -- (#1,#2,#3+#6) -- cycle;
    \filldraw[aresta, fill=#7!20] (#1,#2,#3+#6) -- (#1+#4,#2,#3+#6) -- (#1+#4,#2+#5,#3+#6) -- (#1,#2+#5,#3+#6) -- cycle;}
  % a abraçadeira, sob a base
  \filldraw[aresta, fill=nro-cinza!30] (4.2,4,-1.3) -- (7.8,4,-1.3) -- (7.8,4,0) -- (4.2,4,0) -- cycle;
  \draw[aresta, fill=white] (6,4,-1.9) circle (1.9mm);
  \node[anchor=west, align=left] at (13.2,4,-1.2) {abraçadeira para tubo de 35\,mm,\\braço de 8\,mm e calço de TPU};
  % a base (12 x 9 x 1,8)
  \caixa{0}{0}{0}{12}{9}{1.8}{nro-fundo}
  \foreach \x/\y in {1/1,11/1,1/8,11/8} \fill[nro-synapse!80!black] (\x,\y,1.8) circle (0.9mm);
  \node[anchor=west, align=left] at (13.2,4.5,0.9) {base em PETG, paredes de 2,4\,mm,\\insertos de latão M3};
  % as linhas de montagem (antes da placa e da tampa, que as escondem)
  \foreach \x/\y in {1/1,11/1,1/8,11/8} \draw[nro-claro, densely dashed] (\x,\y,1.8) -- (\x,\y,7.8);
  % a placa (10 x 7,5 x 0,2), com os componentes altos
  \caixa{1}{0.75}{4.2}{10}{7.5}{0.25}{nro-axon}
  \caixa{2}{2}{4.45}{2.5}{2.5}{0.6}{nro-cinza}
  \caixa{6}{1.2}{4.45}{4}{1.2}{0.7}{nro-tinta}
  \caixa{9.6}{3}{4.45}{1.2}{3}{0.9}{nro-synapse}
  \node[anchor=west, align=left] at (13.2,4.5,4.6) {placa de potência, revisão B\\(Atividade 1)};
  % a tampa, com as ranhuras em labirinto
  \caixa{0}{0}{7.8}{12}{9}{1.6}{nro-fundo}
  \foreach \k in {0,...,5} \filldraw[draw=nro-tinta, line width=0.3pt, fill=nro-tinta!55]
    (2+\k*1.3,5,9.4) -- (2.6+\k*1.3,5,9.4) -- (2.6+\k*1.3,8,9.4) -- (2+\k*1.3,8,9.4) -- cycle;
  \draw[aresta, fill=nro-pulso!80] (3.2,2.2,9.4) circle (2.2mm);
  \node[anchor=west, align=left] at (13.2,4.5,8.7) {tampa em PETG com ranhuras em\\labirinto e o botão de parada};
\end{tikzpicture}

  \caption{Vista explodida do invólucro: a base com a abraçadeira, a placa e a tampa.}
  \fonte{Elaborado pela autora.}
  \label{fig:ex-involucro}
\end{figure}

\textbf{A temperatura.} A placa dissipa \SI{2,1}{\watt} dentro do invólucro. Em convecção natural,
com $h \approx \SI{6}{\watt\per\square\meter\per\kelvin}$ e a área externa de
$A = \SI{0,034}{\square\meter}$, a elevação de temperatura estimada é
\begin{equation}
  \Delta T = \frac{P}{h A} = \frac{\SI{2,1}{\watt}}{\SI{6}{\watt\per\square\meter\per\kelvin} \times \SI{0,034}{\square\meter}} \approx \SI{10}{\kelvin},
  \label{eq:ex-termico}
\end{equation}
o que dá \SI{45}{\celsius} com \SI{35}{\celsius} de ambiente, dentro do requisito.

\textbf{A abraçadeira.} Tomei a massa de \SI{0,38}{\kilo\gram} sob 5\,g de vibração
($F = m a \approx \SI{19}{\newton}$) e dimensionei para $F = \SI{50}{\newton}$. O braço da abraçadeira
é uma viga em balanço de $L = \SI{30}{\milli\meter}$ e seção $b \times h$, com a tensão máxima
\begin{equation}
  \sigma = \frac{M c}{I} = \frac{F L \cdot h/2}{b h^3/12} = \frac{6 F L}{b h^2}.
  \label{eq:ex-viga}
\end{equation}
Com $b = \SI{12}{\milli\meter}$ e $h = \SI{5}{\milli\meter}$, a primeira versão dava
$\sigma = \SI{30}{\mega\pascal}$ --- fator de segurança 1,0 contra os \SI{30}{\mega\pascal} que
estimei para o PETG impresso na direção entre camadas (60\,\% dos \SI{50}{\mega\pascal} do material).
Aumentei para $h = \SI{8}{\milli\meter}$: $\sigma = \SI{11,7}{\mega\pascal}$ e fator 2,6.
\end{desenvolvimento}

\begin{resultados}
O conjunto com a placa e a bateria pesou \SI{352}{\gram}. No ensaio térmico, em estimulação contínua
na amplitude máxima, a primeira impressão, sem ranhuras, passou de \SI{49}{\celsius} com
\SI{26}{\celsius} de ambiente; com as ranhuras em labirinto, a temperatura estabilizou em
\SI{38}{\celsius}, uma elevação de \SI{12,5}{\kelvin}, perto dos \SI{10}{\kelvin} estimados em
\eqref{eq:ex-termico} (Figura~\ref{fig:ex-termico}). Extrapolada para \SI{35}{\celsius} de ambiente,
fica em \SI{48}{\celsius}, abaixo dos \SI{60}{\celsius} do requisito.

\begin{figure}[htbp]
  \centering
  % Temperatura do ar dentro do invólucro em estimulação contínua na
% amplitude máxima, a 26 °C de ambiente (exemplo fictício).
\begin{tikzpicture}
  \begin{axis}[
      width=0.86\linewidth, height=58mm,
      xlabel={Tempo (min)}, ylabel={Temperatura interna (\si{\celsius})},
      xmin=0, xmax=60, ymin=20, ymax=65,
      grid=major, grid style={nro-fundo},
      tick label style={font=\scriptsize}, label style={font=\footnotesize},
      legend pos=north west, legend cell align=left,
      legend style={font=\scriptsize, draw=nro-linha, fill=white, fill opacity=0.9, text opacity=1}]
    \addplot[only marks, mark=square*, mark size=1.4pt, nro-pulso-texto]
      coordinates {(0,25.8) (5,33.2) (10,38.1) (15,41.9) (20,44.3) (25,45.6) (30,46.8) (35,48.3) (40,48.4) (45,48.8) (50,49.7) (55,49.5) (60,49.9)};
    \addlegendentry{sem ranhuras (1.ª impressão)}
    \addplot[only marks, mark=*, mark size=1.5pt, nro-axon]
      coordinates {(0,26.1) (5,29.9) (10,32.6) (15,34.6) (20,35.7) (25,36.7) (30,36.7) (35,37.4) (40,38.1) (45,38.1) (50,38.5) (55,37.9) (60,38.3)};
    \addlegendentry{com ranhuras em labirinto}
    \addplot[domain=0:60, samples=61, nro-axon, thin] {26+12.5*(1-exp(-x/14))};
    \addlegendentry{ajuste: $\Delta T_\infty = 12{,}5$\,K, $\tau = 14$\,min}
    \addplot[domain=0:60, nro-cinza, densely dashed] {60};
    \addlegendentry{limite de projeto (60\,°C)}
  \end{axis}
\end{tikzpicture}

  \caption{Temperatura do ar dentro do invólucro em estimulação contínua na amplitude máxima, a
    \SI{26}{\celsius} de ambiente.}
  \fonte{Relatório de execução \texttt{NRO-PRO-003-902}.}
  \label{fig:ex-termico}
\end{figure}

No ensaio de vibração --- três horas na bicicleta, na bancada de rolo ---, nada se soltou, e a
remoção do conjunto levou \SI{40}{\second}. Dois problemas apareceram: os cantos da primeira base
empenaram na impressão, resolvido com borda de adesão e a impressora fechada; e o encaixe da tampa
ficou apertado demais, resolvido ao passar a folga de \num{0,2} para \SI{0,3}{\milli\meter}.
\end{resultados}

\begin{evidencias}
\evidencia{Arquivo de projeto}{involucro-rev1 (FreeCAD e STL)}{O modelo, os desenhos e os arquivos de impressão}{Drive da equipe (acesso a pedido)}
\evidencia{Arquivo controlado}{NRO-PRO-014-902}{O design record do invólucro, com a matriz de decisão e os cálculos}{SOMA › Arquivos}
\evidencia{Arquivo controlado}{NRO-PRO-003-902}{O relatório dos ensaios térmico e de vibração}{SOMA › Arquivos}
\evidencia{Atividade}{NBL-127}{O cartão do invólucro, atribuído a mim e concluído}{SOMA › Atividades › NEBULA}
\end{evidencias}

\begin{aprendizados}
Uma estudante de Engenharia Elétrica projetando uma peça mecânica: usei o Desenho Técnico mais do
que esperava, e a conta de convecção de Fenômenos de Transporte decidiu as ranhuras. O que aprendi fora
da sala: a peça impressa não é o material do catálogo --- a anisotropia entre camadas mudou o
dimensionamento ---, e um requisito de segurança (nenhuma parte de alta tensão acessível) vira
geometria, não só texto. Da próxima vez, eu imprimiria um corpo de prova para medir a resistência entre
camadas, em vez de estimá-la.
\end{aprendizados}
\end{atividade}

%% EXEMPLO FICTÍCIO — Atividade 4: ensaio experimental e validação
\begin{atividade}{Ensaios de bancada da revisão B e sessões de validação com o atleta-piloto}{
  tipo         = ensaio,
  modalidade   = EXT,
  alternativa  = PES,
  horas        = 70,
  inicio       = 2024-09,
  fim          = 2025-02,
  projeto      = {NEBULA, estimulação elétrica funcional para ciclismo},
  papel        = {Responsável pela parte técnica},
  supervisor   = {Rafael Fictício, supervisor do NEBULA; Dra. Carla Fictícia, fisioterapeuta do Depto. Clínico},
  registros    = {NRO-PRO-005-901; NRO-PRO-003-901; NRO-PRO-003-903; NBL-135},
  competencias = {I, II, IV, VI, VII},
  disciplinas  = {Laboratório de Eletrônica; Instrumentação Eletrônica; Probabilidade e Estatística},
}

\begin{contexto}
Uma placa que funciona na bancada ainda não é um sistema que serve a quem vai usá-lo. O plano de
validação do NEBULA (\texttt{NRO-PRO-005-901}) pedia duas etapas depois da montagem: a verificação
da revisão B contra os requisitos elétricos da USRS, na bancada; e doze sessões de treino com o
atleta-piloto --- uma pessoa com lesão medular, voluntária, de fora da UFMG ---, para confirmar que o
sistema faz o que ela precisa: pedalar por mais tempo. As sessões são a parte da ação de extensão da
equipe em que a tecnologia volta para quem motivou o projeto.
\end{contexto}

\begin{contribuicao}
Escrevi as seções de bancada do plano de validação, conduzi os ensaios de bancada e respondi pela
parte técnica das doze sessões: antes de cada uma, conferi a impedância dos eletrodos e carreguei os
parâmetros prescritos; durante, registrei os dados e monitorei a segurança elétrica. Quem conduziu as
sessões, definiu os parâmetros de estimulação e esteve com o atleta foi a fisioterapeuta do
Departamento Clínico; eu não posicionei eletrodos nem ajustei nada sem a indicação dela.
\end{contribuicao}

\begin{desenvolvimento}
\fichatecnica{
  protocolo    = {Plano de validação \texttt{NRO-PRO-005-901}, seções 4 (bancada) e 6 (sessões)},
  instrumentos = {Osciloscópio de 100\,MHz e sonda de corrente de 50\,MHz (calibrados em 2024); câmera térmica; cargas resistivas de 0,1\,\%; multímetro de 6\,½ dígitos},
  local        = {LABBIO, Escola de Engenharia da UFMG},
  amostra      = {Duas placas; um atleta-piloto voluntário, em 12 sessões},
  etica        = {Protocolo aprovado pelo Comitê de Ética em Pesquisa da UFMG, CAAE 00000000.0.0000.0000 (exemplo); termo de consentimento assinado},
  relatorio    = {\texttt{NRO-PRO-003-901} (bancada) e \texttt{NRO-PRO-003-903} (sessões)},
}

\textbf{A bancada.} Cada canal das duas placas foi ensaiado em três cargas de 0,1\,\% (\SI{500}{\ohm},
\SI{1}{\kilo\ohm} e \SI{1,5}{\kilo\ohm}), em seis amplitudes, com cinco repetições
(Figura~\ref{fig:ex-bancada}). A corrente foi medida com a sonda, e a carga de cada fase, pela integral
do pulso no próprio osciloscópio; o desbalanço de carga é a diferença entre as fases dividida pela
carga de uma delas. A incerteza combinada da amplitude, somando a da sonda (1\,\%) e a do canal do
osciloscópio (1,5\,\%) em quadratura, é de 1,8\,\%.

\begin{figure}[htbp]
  \centering
  \resizebox{0.92\linewidth}{!}{% A montagem dos ensaios de bancada (exemplo fictício).
\begin{tikzpicture}[
    bloco/.style={draw=nro-tinta, line width=0.5pt, fill=white, rounded corners=1.2pt,
      minimum height=12mm, text width=31mm, align=center, font=\footnotesize, inner sep=2pt},
    seta/.style={-{Stealth[length=1.8mm]}, line width=0.55pt, draw=nro-tinta},
    x=1mm, y=1mm]
  \node[bloco, fill=nro-synapse!18] (est) at (0,20) {Estimulador\\{\scriptsize placa rev.\,B no invólucro}};
  \node[bloco] (carga) at (52,20) {Carga de teste\\{\scriptsize 500\,$\Omega$, 1\,k$\Omega$ e 1,5\,k$\Omega$ (0,1\,\%)}};
  \node[bloco] (osc) at (104,20) {Osciloscópio\\{\scriptsize 100\,MHz, 4 canais}};
  \node[bloco] (note) at (0,0) {Notebook\\{\scriptsize registro dos parâmetros}};
  \node[bloco] (cam) at (52,0) {Câmera térmica\\{\scriptsize estágio de saída}};
  \node[bloco] (mult) at (104,0) {Multímetro de 6\,½ dígitos\\{\scriptsize calibração de $R_S$}};
  \draw[seta] (est) -- node[above, font=\scriptsize] {canal 1 a 4} (carga);
  \draw[seta] (carga) -- node[above, font=\scriptsize, align=center] {sonda de\\corrente} (osc);
  \draw[seta] (note) -- node[left, font=\scriptsize] {USB} (est);
  \draw[seta, densely dashed] (cam) -- (est.south east);
  \draw[seta, densely dashed] (mult) -- (carga.south east);
\end{tikzpicture}
}
  \caption{A montagem dos ensaios de bancada. As linhas tracejadas são as medições auxiliares.}
  \fonte{Elaborado pela autora.}
  \label{fig:ex-bancada}
\end{figure}

\textbf{As sessões.} Doze sessões de cerca de uma hora, em oito semanas, no LABBIO. Antes de cada uma,
eu media a impedância dos eletrodos e, acima de \SI{1}{\kilo\ohm} --- o limite que saiu da
Atividade~2 ---, a fisioterapeuta trocava o eletrodo. O indicador de desempenho definido no plano foi o
tempo de pedalada contínua até a cadência cair 20\,\% abaixo da inicial. As quatro primeiras sessões
usaram o perfil de estímulo da revisão A; as oito seguintes, o perfil ajustado pela fisioterapeuta com
a faixa maior de largura de pulso que a revisão B permite. Os dados foram registrados sem nenhuma
informação que identifique o atleta.
\end{desenvolvimento}

\begin{resultados}
Na bancada, as duas placas atenderam a todos os requisitos (Tabela~\ref{tab:ex-bancada}). Na carga de
\SI{1,5}{\kilo\ohm}, a corrente ficou limitada a \SI{78}{\milli\ampere}, como a tensão de conformidade
prevê (\SI{118}{\volt}/\SI{1,5}{\kilo\ohm}): nessa carga, as amplitudes acima disso ficam fora da
faixa de operação, e o firmware da revisão C passou a bloqueá-las.

\begin{table}[htbp]
  \centering\small
  \caption{Verificação da revisão B na bancada (duas placas, quatro canais, cinco repetições).}
  \label{tab:ex-bancada}
  \begin{tabular}{|l|c|c|}\hline
    \rowcolor{nro-fundo}\rotulo{Requisito} & \rotulo{Critério de aceite} & \rotulo{Pior resultado} \\ \hline
    Exatidão da amplitude (10 a 100\,mA, até 1\,k$\Omega$) & $\le$\,2\,\% & 0,9\,\% \\ \hline
    Desbalanço de carga entre as fases & $\le$\,2\,\% & 0,8\,\% \\ \hline
    Largura do pulso (50 a 500\,µs) & $\le$\,2\,µs de erro & 0,6\,µs \\ \hline
    Frequência (20 a 50\,Hz) & $\le$\,0,5\,\% & 0,1\,\% \\ \hline
    Disparo do limite de corrente & 105 a 115\,mA & 108 a 111\,mA \\ \hline
  \end{tabular}
  \fonte{Relatório de execução \texttt{NRO-PRO-003-901}.}
\end{table}

Nas sessões, o tempo até a queda de 20\,\% da cadência passou de 11 minutos, em média, nas quatro
primeiras sessões, para 19 minutos nas quatro últimas. Não houve evento adverso. O limite de corrente
por hardware disparou duas vezes, ambas por eletrodo descolado durante a pedalada: a saída foi
desligada antes de qualquer desconforto, como o projeto prevê. O atleta pediu que a parada de
emergência ficasse ao alcance da mão direita --- mudança feita no invólucro.
\end{resultados}

\begin{evidencias}
\evidencia{Arquivo controlado}{NRO-PRO-005-901}{O plano de validação, com as seções de bancada que escrevi}{SOMA › Arquivos}
\evidencia{Arquivo controlado}{NRO-PRO-003-901}{O relatório dos ensaios de bancada da revisão B}{SOMA › Arquivos}
\evidencia{Arquivo controlado}{NRO-PRO-003-903}{O relatório das doze sessões, sem dados pessoais}{SOMA › Arquivos (classe confidencial)}
\evidencia{Atividade}{NBL-135}{O cartão da validação, atribuído a mim e concluído}{SOMA › Atividades › NEBULA}
\evidencia{Presença}{Check-ins de set./2024 a fev./2025}{A presença no LABBIO nos dias de ensaio e de sessão}{SOMA › Equipe › Presença (a pedido)}
\end{evidencias}

\begin{aprendizados}
A bancada pôs em prática o que Instrumentação Eletrônica ensina sobre incerteza, e o que Probabilidade e
Estatística ensina sobre repetição; o Laboratório de Eletrônica me deu o hábito do procedimento
escrito. O que só as sessões ensinaram: trabalhar ao lado de uma profissional de saúde, respeitando o
que é dela; que um protocolo aprovado por comitê de ética é compromisso, não burocracia; e que o
requisito mais importante das doze sessões veio do próprio atleta --- a posição do botão de parada.
\end{aprendizados}
\end{atividade}

%% EXEMPLO FICTÍCIO — Atividade 5: ensino e formação de membros
\begin{atividade}{Treinamento de segurança elétrica em bancada, para os trainees de 2025/1}{
  tipo         = ensino,
  modalidade   = ENS,
  alternativa  = VPC,
  horas        = 20,
  inicio       = 2025-03,
  fim          = 2025-04,
  projeto      = {Formação de membros, com o Depto. de Pessoal},
  papel        = {Autora e instrutora},
  supervisor   = {Mariana Fictícia, gerente do Depto. de Pessoal},
  registros    = {NRO-TRE-905; PES-58},
  competencias = {V, VI, VII},
  disciplinas  = {Laboratório de Eletrônica; Instalações Elétricas},
}

\begin{contexto}
Com a revisão B, a bancada do NEBULA passou a ter \SI{120}{\volt} contínuos numa placa aberta. O
treinamento de segurança que a equipe tinha era um texto, de 2022, sobre a rede elétrica; não tratava
de tensão contínua, de capacitores carregados nem de fontes com limite de corrente. Em 2024, um membro
levou um choque leve do boost da revisão A --- sem consequência, mas registrado como ocorrência. Com
catorze trainees chegando em 2025/1, a equipe precisava de um treinamento prático, com verificação.
\end{contexto}

\begin{contribuicao}
Escrevi os três módulos no SOMA, as dez questões da verificação de conhecimento e o checklist de
energização; gravei dois vídeos curtos; e dei as duas sessões práticas de uma hora e meia. O supervisor
do projeto revisou o conteúdo técnico, e a gerente do Depto. de Pessoal atribuiu o treinamento aos
trainees como obrigatório.
\end{contribuicao}

\begin{desenvolvimento}
\fichatecnica{
  formato       = {Treinamento no SOMA (três módulos, com verificação) e duas sessões práticas},
  treinamentos  = {\texttt{NRO-TRE-905} Rev.~A e Rev.~B},
  publico       = {Os 14 trainees de 2025/1 e 3 membros do NEBULA},
  cargahoraria  = {3\,h (duas sessões de 1,5\,h)},
  participantes = {17},
  material      = {3 módulos, 2 vídeos, 10 questões e o checklist de energização de placas de alta tensão},
}

O treinamento foi escrito para quem nunca pôs a mão numa bancada de alta tensão
(Tabela~\ref{tab:ex-modulos}). Cada módulo termina com questões, e a aprovação pede 70\,\% --- a nota
mínima padrão do SOMA. A sessão prática só começa depois da aprovação.

\begin{table}[htbp]
  \centering\small
  \caption{Os módulos do treinamento \texttt{NRO-TRE-905}.}
  \label{tab:ex-modulos}
  \begin{tabular}{|c|p{47mm}|p{85mm}|}\hline
    \rowcolor{nro-fundo}\rotulo{Módulo} & \rotulo{Tema} & \rotulo{O que ensina} \\ \hline
    1 & Por que 120\,V contínuos importam & A corrente no corpo e os seus efeitos; por que a tensão contínua prende; a energia guardada nos capacitores. \\ \hline
    2 & O procedimento de bancada & A fonte com limite de corrente; a descarga dos capacitores antes de tocar; uma mão só na bancada; o que fazer num acidente. \\ \hline
    3 & A prática & A energização da placa revisão B com o checklist, passo a passo, com a instrutora ao lado. \\ \hline
  \end{tabular}
  \fonte{Elaborado pela autora.}
\end{table}
\end{desenvolvimento}

\begin{resultados}
Os dezessete participantes concluíram o treinamento, com nota média de 87\,\% na verificação (a menor
foi 70\,\%). Duas questões tiveram mais de metade de erros e, lendo as respostas, vi que o problema era
a redação, não o conteúdo: reescrevi-as na revisão B do treinamento, publicada em abril de 2025. Desde
então, nenhuma ocorrência de segurança elétrica foi registrada na equipe, e o checklist passou a ficar
impresso na bancada.
\end{resultados}

\begin{evidencias}
\evidencia{Treinamento}{NRO-TRE-905 Rev. B}{O treinamento de minha autoria, publicado no SOMA}{SOMA › Treinamentos}
\evidencia{Relatório}{Acompanhamento do NRO-TRE-905}{Quem concluiu, quando e com que nota}{SOMA › Treinamentos (a pedido)}
\evidencia{Atividade}{PES-58}{O cartão do treinamento no quadro do Depto. de Pessoal, concluído}{SOMA › Atividades}
\end{evidencias}

\begin{aprendizados}
Ensinar me obrigou a saber de verdade o que eu fazia por hábito na bancada. O Laboratório de Eletrônica
e Instalações Elétricas me deram a base; o que aprendi ensinando foi escrever para quem não sabe, e
medir se a pessoa aprendeu --- as duas questões mal escritas me ensinaram mais do que as oito boas.
\end{aprendizados}
\end{atividade}

%% EXEMPLO FICTÍCIO — Atividade 6: evento, competição ou visita técnica
\begin{atividade}{Competição internacional de ciclismo com estimulação elétrica funcional}{
  tipo         = evento,
  modalidade   = EVT,
  alternativa  = COMP,
  horas        = 60,
  inicio       = 2024-10,
  fim          = 2024-10,
  projeto      = {NEBULA, estimulação elétrica funcional para ciclismo},
  papel        = {Técnica responsável pelo estimulador},
  supervisor   = {Rafael Fictício, supervisor do projeto NEBULA},
  registros    = {EXT-901; NRO-DIR-006-901},
  competencias = {IV, V, VI},
}

\begin{contexto}
A competição reúne equipes de vários países em provas com tecnologias assistivas; na prova de ciclismo
com estimulação elétrica funcional, o atleta pedala uma distância fixa, e o sistema é avaliado também
por uma banca técnica. Para o NEBULA, é o teste mais duro que existe: o sistema tem de funcionar longe
do laboratório, depois de uma viagem, sob o regulamento de outra instituição.
\end{contexto}

\begin{contribuicao}
Fui a técnica responsável pelo estimulador: montei o checklist de preparação, levei uma placa reserva
configurada, fiz o plano de carga das baterias, operei o estimulador nas provas e apresentei a parte
eletrônica do sistema à banca técnica, em inglês.
\end{contribuicao}

\begin{desenvolvimento}
\fichatecnica{
  evento     = {Competição internacional de tecnologias assistivas, prova de ciclismo com FES},
  local      = {Europa},
  datas      = {Quatro dias em outubro de 2024},
  funcao     = {Técnica responsável pelo estimulador},
  resultado  = {4.º lugar entre 9 equipes},
  declaracao = {\texttt{NRO-DIR-006-901}},
}

\textbf{A preparação} (20 horas, em setembro). O checklist de preparação cobriu a placa reserva, os
eletrodos, os cabos e o transporte das baterias de lítio, que as companhias aéreas só aceitam abaixo de
\SI{100}{\watt\hour} por unidade e com a documentação do fabricante --- as nossas têm
\SI{29}{\watt\hour}. Simulei, com o atleta e a fisioterapeuta, duas provas completas no LABBIO.

\textbf{A competição} (quatro dias de cerca de dez horas). Conferência dos equipamentos na chegada,
treinos livres, a apresentação à banca técnica e as provas. Em cada prova, eu conferia a impedância dos
eletrodos, carregava o perfil de estímulo e acompanhava a telemetria.
\end{desenvolvimento}

\begin{resultados}
A equipe terminou em 4.º lugar entre 9. O estimulador funcionou em todas as provas; um eletrodo descolou
na segunda, e o limite de corrente por hardware desligou o canal como devia --- o atleta completou a
prova com três canais. Da conversa com as outras equipes, trouxemos a ideia do controle da cadência em
malha fechada, que virou requisito do NEBULA para 2025.
\end{resultados}

\begin{evidencias}
\evidencia{Declaração}{NRO-DIR-006-901}{A declaração de participação, com a função e as horas}{auth.neurodynamics.dev (código verificador)}
\evidencia{Registro}{EXT-901}{O registro do evento no SOMA, aprovado por duas pessoas}{SOMA › Serviços › Eventos}
\evidencia{Resultado}{Classificação oficial}{A colocação da equipe na prova}{Site da competição}
\end{evidencias}

\begin{aprendizados}
Operar um sistema meu diante de uma banca de outro país ensinou mais sobre preparação do que qualquer
teste de bancada: o checklist salvou a segunda prova. E ver sistemas de outras equipes mostrou o que o
nosso ainda não fazia.
\end{aprendizados}
\end{atividade}

%% EXEMPLO FICTÍCIO — Atividade 7: gestão e liderança
\begin{atividade}{Supervisão do subsistema eletrônico do NEBULA}{
  tipo         = gestao,
  modalidade   = GES,
  alternativa  = VPC,
  horas        = 60,
  inicio       = 2025-01,
  fim          = 2025-12,
  projeto      = {NEBULA, estimulação elétrica funcional para ciclismo},
  papel        = {Supervisora do subsistema},
  supervisor   = {Rafael Fictício, supervisor do projeto NEBULA},
  registros    = {OKRs 2025/1 e 2025/2 do NEBULA; NRO-PRO-013-903; NBL-150 a NBL-189},
  competencias = {IV, VI, VIII},
}

\begin{contexto}
Em janeiro de 2025, pela política de progressão funcional (\texttt{NRO-PES-013}), assumi a supervisão
do subsistema eletrônico do NEBULA. O momento: a revisão B funcionava, mas a errata tinha sete itens e
o controle da cadência em malha fechada, trazido da competição, disputava o mesmo semestre. O time
eram duas pessoas da equipe e dois trainees recém-chegados.
\end{contexto}

\begin{contribuicao}
Planejei os dois semestres, conduzi as sprints, decidi as prioridades, instituí a revisão de projeto por
duas pessoas para qualquer mudança de placa, acompanhei a formação dos dois trainees e preparei a minha
sucessão.
\end{contribuicao}

\begin{desenvolvimento}
\fichatecnica{
  cargo       = {Supervisora do subsistema eletrônico, de jan. a dez./2025},
  time        = {4 pessoas: 2 membros e 2 trainees},
  rituais     = {Sprints de duas semanas; acompanhamento diário assíncrono; revisão mensal com o supervisor do projeto; OKRs semestrais},
  documentos  = {\texttt{NRO-PRO-013-903} (a decisão da revisão C); checklist de energização; plano de transição},
  indicadores = {Cartões concluídos no prazo; revisões aprovadas; pessoas formadas; ocorrências de segurança},
}

\textbf{O planejamento.} O objetivo de 2025/1 foi \enquote{liberar a revisão C, validada, até
setembro}, com três resultados-chave: a errata inteira resolvida, a placa aprovada na bancada e a
documentação revisada. O controle em malha fechada ficou para 2026: fazê-lo junto com a revisão C
atrasaria as duas, e registrei a decisão, com as alternativas, em \texttt{NRO-PRO-013-903}.

\textbf{O acompanhamento.} Cada sprint começava com os cartões do quadro distribuídos e terminava com a
demonstração do que ficou pronto. Toda mudança de placa passou a ter dois revisores antes de ir para
fabricação --- a lição do footprint espelhado da revisão B.

\textbf{A sucessão.} De setembro em diante, trabalhei em par com quem me sucedeu: ele conduziu as
sprints, eu revisei. O plano de transição lista o estado de cada frente, os arquivos e os contatos.
\end{desenvolvimento}

\begin{resultados}
\begin{table}[htbp]
  \centering\small
  \caption{Os indicadores da supervisão em 2025.}
  \label{tab:ex-indicadores}
  \begin{tabular}{|l|c|}\hline
    \rowcolor{nro-fundo}\rotulo{Indicador} & \rotulo{Resultado} \\ \hline
    Cartões concluídos no prazo & 31 de 36 (86\,\%) \\ \hline
    Revisão C liberada & 22/9/2025 (a meta era 15/9) \\ \hline
    Trainees efetivados no subsistema & 2 de 2 \\ \hline
    Ocorrências de segurança elétrica & nenhuma \\ \hline
  \end{tabular}
  \fonte{OKRs e quadro de atividades do NEBULA no SOMA.}
\end{table}
A revisão C saiu com uma semana de atraso, porque um componente importado demorou; a errata foi
resolvida por inteiro, e nenhuma placa voltou da revisão de projeto com erro de footprint.
\end{resultados}

\begin{evidencias}
\evidencia{Registro}{OKRs 2025/1 e 2025/2 do NEBULA}{Os objetivos, os resultados-chave e o fechamento de cada semestre}{SOMA › OKRs}
\evidencia{Arquivo controlado}{NRO-PRO-013-903}{A decisão de adiar o controle em malha fechada, com as alternativas}{SOMA › Arquivos}
\evidencia{Atividade}{NBL-150 a NBL-189}{Os cartões do subsistema em 2025, com o histórico de cada um}{SOMA › Atividades › NEBULA}
\evidencia{Registro}{Apontamentos semanais de 2025}{A avaliação de assiduidade e entregas do time}{SOMA › Equipe (a pedido)}
\end{evidencias}

\begin{aprendizados}
Liderar ensinou que decidir o que não fazer é a parte difícil: adiar a malha fechada foi a decisão mais
importante do ano. O curso não tem uma disciplina sobre isso; tem a Engenharia Econômica, que ajudou a
pesar as alternativas, e o hábito de documentar que a equipe me deu desde o trainee.
\end{aprendizados}
\end{atividade}


%% ============================================================
\parte[O conjunto: as horas e a base delas, as competências desenvolvidas, a relação com o curso e a
  autoavaliação pelo mesmo barema que o professor usa.]{Síntese}

\section{As horas pedidas}
\quadrohoras

\subsection{Memória de cálculo}
Mantive uma planilha semanal de horas durante todo o período, conferida com os check-ins no LABBIO e
com as estimativas dos cartões de cada atividade no SOMA. Nos 32 meses de participação, descontado o
afastamento, somei cerca de 1.300 horas --- em média, perto das dez horas semanais declaradas.

Peço 440 horas: só as que se ligam diretamente às sete atividades descritas e que têm evidência. Para
cada atividade, somei as horas dos cartões dela (os códigos estão na ficha de cada uma) e conferi, nos
dias de bancada, com os check-ins. Ficaram de fora o período trainee, as reuniões gerais, o trabalho
em outros subsistemas e toda hora sem cartão. A Atividade~6 conta 20 horas de preparação e quatro dias
de dez horas de competição, como diz a declaração \texttt{NRO-DIR-006-901}.

\section{As competências desenvolvidas}
\matrizcompetencias

\section{A relação com o curso}
\quadrodisciplinas

As atividades aplicaram, com um usuário real, o núcleo do curso de Engenharia Elétrica: a fonte de
corrente com amplificador operacional e o dimensionamento da dissipação (Atividade~1), o transitório de
primeira ordem e a margem de fase (Atividade~2), a incerteza de medição (Atividade~4). E ensinaram o que
o curso não ensina: projetar para a segurança de uma pessoa ligada ao circuito, ler uma norma de
equipamento eletromédico \parencite{iec60601-1,iec60601-2-10}, documentar para que outra pessoa continue
o trabalho, trabalhar com uma profissional de saúde e liderar um time com prazo.

\section{Autoavaliação}
\autoavaliacao{
  aderencia    = 9,
  profundidade = 8,
  contribuicao = 9,
  evidencias   = 9,
  resultados   = 8,
  comunicacao  = 8,
}
\begin{itemize}
  \item \textbf{Aderência (9):} as sete atividades são de engenharia elétrica, e a relação com as
    disciplinas está mostrada em cada uma; a do invólucro é mais mecânica que elétrica.
  \item \textbf{Profundidade (8):} há cálculo, simulação validada e decisões justificadas; faltou medir a
    resistência do PETG entre camadas, que estimei.
  \item \textbf{Contribuição (9):} cada atividade separa o que foi meu do que foi da equipe e nomeia
    quem fez o resto.
  \item \textbf{Evidências (9):} todas as atividades têm registros no SOMA; parte dos arquivos de
    projeto fica em repositórios com acesso a pedido.
  \item \textbf{Resultados (8):} os resultados são medidos e comparados com os requisitos; o indicador
    das sessões tem só um usuário.
  \item \textbf{Comunicação (8):} o texto é direto; algumas figuras são esquemáticas, não os desenhos
    completos.
\end{itemize}

\section{Considerações finais}
A equipe me deu o que o curso não podia dar: a responsabilidade por um sistema que uma pessoa usa. O
que deixei para ela: a revisão B e a revisão C da placa, o invólucro, o treinamento de segurança que os
trainees fazem, o checklist que fica na bancada, os registros de decisão que explicam o porquê de cada
escolha e um sucessor que conduziu o subsistema por três meses antes da minha saída.

\referencias

%% ============================================================
\pareceres

\anexoevidencias
\anexo{Declaração de vínculo emitida pelo SOMA}{anexos/declaracao-vinculo.pdf}
\anexo{Declaração de participação na competição}{anexos/declaracao-participacao.pdf}

\end{document}
